\chapter{Fungal Infections}
\index{Fungal Infections}

\section*{Definition and Overview}
Fungal infections, also called mycoses, are caused by microscopic fungi that live on the skin, in the air, and in soil. Most fungal infections affect the skin, nails, and hair, including tinea (ringworm, athlete's foot, jock itch), tinea versicolor, and nail fungus, and these are common, usually minor, and treatable with antifungal creams or tablets. Less commonly, fungi cause serious infections of the lungs, blood, and internal organs, mainly in people with weakened immune systems, and these require hospital treatment. This chapter covers the common fungal infections, how they spread, and how they are treated, and it highlights when fungal infection is serious.

\section*{Epidemiology}
Skin and nail fungal infections are among the most common conditions seen in general practice. Athlete's foot affects around 1 in 5 people at some point, tinea of the body and scalp is common in children, and fungal nail infection affects around 10% of adults, rising with age. Fungal infections thrive in warm, moist conditions and spread in communal showers, gyms, and swimming pools. Serious invasive fungal infections, including candidiasis of the bloodstream and aspergillosis, are uncommon but increasing, occurring mainly in hospitalised patients, people with cancer, transplant recipients, and others with weakened immunity.

\section*{Aetiology and Pathophysiology}
\begin{itemize}
    \item \textbf{Dermatophytes:} fungi that infect skin, hair, and nails, causing ringworm (tinea corporis), athlete's foot (tinea pedis), jock itch (tinea cruris), and scalp ringworm (tinea capitis)
    \item \textbf{Yeasts:} candida, which causes thrush and skin fold infections, and malassezia, which causes tinea versicolor and dandruff
    \item \textbf{Moulds:} such as aspergillus, which causes allergy and serious lung and blood infection in vulnerable people
    \item \textbf{Transmission:} skin fungi spread by direct contact with infected people, animals, and contaminated surfaces such as shower floors and towels
    \item \textbf{Risk factors for skin infection:} warm, moist skin, sweating, tight clothing, communal facilities, diabetes, and reduced immunity
    \item \textbf{Risk factors for invasive infection:} chemotherapy, transplant, HIV, intensive care, and long courses of steroids
    \item \textbf{Mechanism:} fungi invade the keratin of skin, hair, and nails, or, in invasive disease, grow in the lungs and spread through the blood
    \item \textbf{The immune system:} most people control fungal infections effectively; serious disease reflects immune failure
\end{itemize}

\section*{Clinical Features}
\begin{itemize}
    \item Ringworm: a red, scaly, ring-shaped patch with a clearer centre, often itchy
    \item Athlete's foot: itching, scaling, and cracking between the toes, spreading to the soles
    \item Jock itch: a red, itchy rash in the groin
    \item Scalp ringworm: scaly patches with hair loss, most common in children
    \item Fungal nail infection: thickened, discoloured, crumbling nails
    \item Tinea versicolor: pale or dark, fine, scaly patches on the chest and back
    \item Thrush: white patches in the mouth, or itching and discharge in the vagina
    \item Invasive disease: fever, cough, breathlessness, and signs of sepsis, mainly in immunocompromised people
\end{itemize}

\section*{Differential Diagnosis}
Skin rashes and nail changes have many causes.
\begin{itemize}
    \item Eczema and psoriasis, which can mimic tinea
    \item Contact dermatitis and allergic rashes
    \item Bacterial skin infections, including impetigo
    \item Seborrhoeic dermatitis and dandruff
    \item Psoriasis of the nails, which resembles fungal nail infection
    \item Herpes and other viral rashes
    \item Invasive fungal disease is distinguished by its setting and testing
\end{itemize}

\section*{When to Seek Medical Care}
Most skin fungal infections can be treated with over-the-counter creams, but medical care is needed if the infection is extensive, persistent, on the scalp or nails, in a child, or not responding to self-care. People with weakened immunity who develop fever or respiratory symptoms should be assessed promptly.

\textbf{Contact your local emergency services for:}
\begin{itemize}
    \item Fever, breathlessness, or chest pain in an immunocompromised person
    \item Signs of sepsis: confusion, rapid heart rate, and collapse
    \item Rapidly spreading skin infection with fever
    \item Any emergency symptoms
\end{itemize}

\section*{Diagnosis and Investigations}
\begin{itemize}
    \item \textbf{Clinical examination:} most skin infections are diagnosed on appearance and location
    \item \textbf{Skin scrapings and nail clippings:} examined under the microscope and cultured to confirm the fungus
    \item \textbf{Wood's lamp:} some fungi fluoresce under ultraviolet light
    \item \textbf{Blood tests and cultures:} for invasive disease, including blood cultures and antigen tests
    \item \textbf{Lung imaging:} chest X-ray and CT for suspected lung infection
    \item \textbf{Tissue biopsy:} in uncertain or invasive cases
\end{itemize}

\section*{Management}
\begin{itemize}
    \item \textbf{Athlete's foot and jock itch:} antifungal creams (clotrimazole, terbinafine) for 2--4 weeks, with keeping the area dry
    \item \textbf{Ringworm:} antifungal cream; widespread or scalp infection needs antifungal tablets
    \item \textbf{Fungal nail infection:} antifungal tablets for several months, as creams rarely penetrate the nail
    \item \textbf{Tinea versicolor:} antifungal shampoo or cream
    \item \textbf{Thrush:} antifungal lozenges or creams, and treatment of predisposing factors
    \item \textbf{Invasive fungal infection:} hospital care with intravenous antifungal medicines such as fluconazole or amphotericin, and treatment of the underlying condition
    \item \textbf{Hygiene measures:} keep skin dry, change socks, avoid sharing towels, and treat household members and pets where relevant
    \item \textbf{Prevent reinfection:} treat shoes, floors, and communal surfaces, and continue treatment for the full course
\end{itemize}

\section*{Complications}
\begin{itemize}
    \item Spread of infection to other body areas and to family members
    \item Chronic, recurrent infection without complete treatment
    \item Secondary bacterial infection of broken skin
    \item Permanent nail damage with untreated nail fungus
    \item Invasive fungal disease causing lung damage, organ failure, and death in immunocompromised people
    \item Resistance to antifungal medicines with incomplete courses
\end{itemize}

\section*{Prognosis and Outcomes}
Most fungal infections respond well to treatment. Skin infections clear within weeks with antifungal creams, nail infections respond to longer courses of tablets, and recurrences are prevented with hygiene measures. Invasive fungal infection has a serious outlook in immunocompromised people but responds to prompt treatment in many cases. The key to good outcomes is completing treatment courses, addressing predisposing factors such as diabetes, and recognising when an infection in a vulnerable person requires urgent care.

\section*{Prevention and Self-Care}
\begin{itemize}
    \item Keep skin clean and dry, especially between the toes and in skin folds
    \item Wear shower shoes in communal showers, gyms, and pools
    \item Change socks daily and alternate shoes so they dry out
    \item Do not share towels, shoes, or clothing
    \item Complete the full course of antifungal treatment, even when symptoms improve
    \item Treat pets and household members with ringworm
    \item Control blood sugar in diabetes
    \item Seek early medical care for fever or respiratory symptoms if immunocompromised
\end{itemize}

\section*{Sources and Further Reading}
The following reputable sources provide further information for healthcare students and patients seeking reliable, current advice about fungal infections.
\begin{itemize}
    \item Healthdirect Australia. \textit{Fungal infections.} https://www.healthdirect.gov.au/fungal-infections
    \item The Australasian College of Dermatologists. \textit{Fungal skin infections.} https://www.dermcoll.edu.au/
    \item Better Health Channel. \textit{Fungal infections.} https://www.betterhealth.vic.gov.au/
    \item NHS. \textit{Fungal infections.} https://www.nhs.uk/conditions/fungal-infections/
    \item Centers for Disease Control and Prevention. \textit{Fungal diseases.} https://www.cdc.gov/
    \item World Health Organization. \textit{Fungal infections.} https://www.who.int/
\end{itemize}
